\section{Showcase: a constrained two-degree-of-freedom system with delayed PD control}\label{sec:showcase}

All elements of the method are now demonstrated on a single mechanical
example that is simple enough to verify by hand yet exercises every
nontrivial feature at once: multiple degrees of freedom, a feedback delay,
and an algebraic constraint, i.e., a delay differential--algebraic system
with a singular mass matrix.

\subsection{Model and automatic extraction}\label{sec:showcase-model}

\begin{figure}[tb]
  \centering
  \maybefig{fig_showcase_hybrid}{0.72\linewidth}
  \caption{Stability chart of the constrained two-degree-of-freedom system
  with delayed PD control. Colored background: combined
  field~\eqref{eq:coloring} of the brute-force sweep (smooth spectral gap
  $\sest$ inside the stable domain, integer number of unstable roots
  outside). Black curve: stability boundary traced by MDBM with the scalar
  objective~\eqref{eq:objective}. Red ellipse and star: largest inscribed
  ellipse and the most robust parameter point. The markers S and U denote
  the stable and unstable points analyzed in
  Fig.~\ref{fig:walkthrough}.}\label{fig:showcase}
\end{figure}

The main structure of mass $m_1$ is statically unstable
(inverted-pendulum-like, modeled by the negative net suspension stiffness
$k_1<0$ and damping $c_1$); it carries an absorber unit of mass $m_2$
attached through a spring and damper $(k_2, c_2)$, and an instrumentation
block of mass $m_3$ is \emph{rigidly bolted} to the absorber mass. Instead of
eliminating the rigid connection by hand, the model keeps it as an algebraic
constraint with the internal force $F_{\mathrm c}$ -- exactly the form in
which a simulation model of a larger assembly would be written. A digital
controller measures the motion of the main mass and applies the delayed PD
force $u(t) = -P\,x_1(t-\tau) - D\,\dot x_1(t-\tau)$ to it. With the state
vector $\vect z = (x_1,x_2,x_3,v_1,v_2,v_3,F_{\mathrm c})$, the descriptor
form~\eqref{eq:ddae} reads
\begin{equation}\label{eq:showcase}
  \underbrace{\operatorname{diag}(1,1,1,m_1,m_2,m_3,0)}_{\textstyle\mat E}\,
  \dot{\vect z}(t) =
  \begin{pmatrix}
    v_1\\[1pt] v_2\\[1pt] v_3\\[1pt]
    -k_1 x_1 - c_1 v_1 + k_2(x_2-x_1) + c_2(v_2-v_1) + u(t)\\[1pt]
    -k_2(x_2-x_1) - c_2(v_2-v_1) + F_{\mathrm c}\\[1pt]
    -F_{\mathrm c}\\[1pt]
    x_3 - x_2
  \end{pmatrix},
\end{equation}
with the numerical values $m_1=1$, $m_2=0.3$, $m_3=0.2$, $k_1=-1$, $k_2=1$,
$c_1=c_2=0.05$ and $\tau=0.5$. The last row of~\eqref{eq:showcase} carries no
time derivative -- the mass matrix is singular -- and the model is passed to
the extraction of Section~\ref{sec:extraction} exactly as written, in its
simulation (right-hand-side) form. For verification, the constraint is also
eliminated by hand ($x_3=x_2$, merged mass $m_{23}=m_2+m_3$), which yields
the reduced quasi-polynomial
\begin{equation}\label{eq:showcasered}
  \Dfun_{\mathrm{red}}(\lam)=\det\!\begin{pmatrix}
   m_1\lam^2+(c_1{+}c_2)\lam+k_1{+}k_2+(P{+}D\lam)\ee^{-\lam\tau} &
   -(c_2\lam+k_2)\\
   -(c_2\lam+k_2) &
   m_{23}\lam^2+c_2\lam+k_2
  \end{pmatrix}.
\end{equation}
The ratio of the automatically extracted $7\times7$ determinant to
$\Dfun_{\mathrm{red}}$ is constant in $\lam$ to machine precision
(deviation below $10^{-15}$ in the computed charts), which validates the
mass-matrix extraction path end to end; the estimated leading
order~\eqref{eq:npow} is $n\approx4.000$, as expected for two mechanical
degrees of freedom. Having verified the two forms identical, the results
below use the cheaper hand-derived form~\eqref{eq:showcasered}; what the
automatic extraction costs instead is measured in
Section~\ref{sec:extraction-cost}.

\subsection{The chart and the phase integration at work}\label{sec:showcase-chart}

\begin{figure}[tb]
  \centering
  \maybefig{fig_method_walkthrough}{\linewidth}
  \caption{The method at the stable point S and unstable point U of
  Fig.~\ref{fig:showcase}. (a)~Normalized locus of
  $\Dfun(\ii\omega)$: the stable curve avoids the origin, the unstable one
  encircles it. (b)~Phase slope $\theta'(\omega)$ with the accepted steps of
  the adaptive solver (dots): steps cluster automatically at the
  near-singular peaks that indicate close characteristic roots.
  (c)~Accumulated phase up to $\wmax=10^{4}$ on a logarithmic frequency
  axis: past the resonances the solver keeps resolving the slowly decaying
  ($\omega^{-1}$) delay ripple of this velocity-feedback system with
  moderate, growing steps.}\label{fig:walkthrough}
\end{figure}

Figure~\ref{fig:showcase} shows the resulting stability chart over the
control-gain plane $(P,D)$. Since the uncontrolled structure is unstable,
the stable domain is a bounded island: its left boundary is the static
(divergence) limit, while the remaining boundaries are Hopf-type curves whose
crossing frequencies are delivered by $\west$. The combined
coloring~\eqref{eq:coloring} displays the interior spectral gap -- deep blue
marks the fastest-decaying configurations -- and the integer plateaus
$\Zint=1,2,3,4$ outside, all in one map. The black MDBM curve is the
high-resolution stability boundary traced from the scalar
objective~\eqref{eq:objective}, and the red star marks the most robust
gain combination found by the largest inscribed ellipse.

Figure~\ref{fig:walkthrough} dissects the single-point computation at the
two marked points. At the stable point S the locus of $\Dfun(\ii\omega)$
avoids the origin and the accumulated phase settles at $\Delta\theta=n\pi/2$,
giving $\Zraw\approx0$; at the unstable point U the locus encircles the
origin and $\Zraw\approx2$. Panel~(b) makes the central numerical argument
of this paper visible: the accepted steps of the adaptive solver concentrate
at the sharp peaks of $\theta'$ near the lightly damped resonances --
precisely the features a fixed frequency grid tends to miss. Panel~(c) shows
the same march continuing over more than two further decades of frequency:
the delayed \emph{velocity} feedback of this system produces a phase ripple
decaying only like $\omega^{-1}$ (Section~\ref{sec:phaseode}), so the solver
keeps taking moderate steps through the tail. The chart studies below use
the standard window $\wmax=10^{4}$ -- its truncated tail,
$\approx0.2/\wmax\approx2\cdot10^{-5}$, is far below the rounding threshold,
and its role in the error budget is measured explicitly -- while the
single-point convergence study of Section~\ref{sec:performance-convergence}
removes even that floor by integrating to $\wmax=10^{6}$. Either way the
truncation is exposed by the integer-residual diagnostic, never hidden.

\subsection{Tolerance as a universal accuracy dial}\label{sec:showcase-tolerance}

\begin{figure}[tb]
  \centering
  \maybefig{fig_tolerance_charts}{\linewidth}
  \caption{The showcase chart computed with solver tolerances
  $10^{-2}$, $10^{-4}$ and $10^{-6}$ (left to right); the
  wall-clock time of each full sweep is given in the panel titles. The charts
  are visually near-identical even at $10^{-2}$, where the count is already
  uncertain: the tolerance buys certainty rather than appearance, cf.\
  Fig.~\ref{fig:errorfield} and Table~\ref{tab:tolerance}. The study uses
  the standard chart window $\wmax=10^{4}$ and the reduced
  form~\eqref{eq:showcasered} of the characteristic
  function.}\label{fig:tolerance}
\end{figure}

\begin{figure}[tb]
  \centering
  \maybefig{fig_error_field}{\linewidth}
  \caption{Self-validating error: fields and histograms of the integer
  residual $\varepsilon=|\Zraw-\operatorname{round}(\Zraw)|$,
  cf.~\eqref{eq:eps}, for the three tolerances of
  Fig.~\ref{fig:tolerance}. Dashed red line: the requested tolerance; dotted
  black line: the practical rule $\varepsilon\approx100\,\mathrm{tol}$. The
  residual scales with the tolerance uniformly over the parameter plane,
  amplified by the accepted-step count of the frequency march, until the
  tail truncated at $\wmax=10^{4}$ (of order $2\cdot10^{-5}$) becomes the
  dominant error.}\label{fig:errorfield}
\end{figure}

\begin{table}[tb]
  \centering
  \caption{Tolerance study on the showcase system
  (grid as in Fig.~\ref{fig:tolerance}): sweep time, integer residual
  statistics, the number of grid points where the rounded count would be
  ambiguous ($\varepsilon\ge1/4$), the number whose $\Zint$ actually deviates
  from the $10^{-6}$-tolerance reference, and the subset of those that flip
  the stable/unstable classification. The ``uncertain'' and ``wrong'' columns
  measure genuinely different things: a skipped peak shifts $\Zraw$ by
  exactly one and leaves it just as close to an integer, so a point can be
  wrong while looking certain -- which is precisely why the
  cross-check~\eqref{eq:crosscheck} is needed.}\label{tab:tolerance}
  \maybeinput{tables/tab_tolerance.tex}
\end{table}

Figure~\ref{fig:tolerance} shows the identical chart computed with
tolerances $10^{-2}$, $10^{-4}$ and $10^{-6}$, together with the
measured sweep times, and Fig.~\ref{fig:errorfield} shows the corresponding
integer-residual fields~\eqref{eq:eps}; the ladder deliberately starts
beyond the safe range. Three observations matter for practice. First, the
tolerance buys \emph{certainty}, not appearance. The panels are visually
near-identical, yet the loosest one is genuinely broken: measured against
the $10^{-6}$ reference, $10^{-2}$ miscounts $349$ of the $8000$ points and
flips the stable/unstable classification on $89$
(Table~\ref{tab:tolerance}). That failure is anything but silent -- the
residual field of Fig.~\ref{fig:errorfield} lights up across the whole
chart, with $210$ points beyond the rounding threshold -- and one decade
into the safe range the problem is essentially gone: at $10^{-4}$ a single
count of the $8000$ deviates, by one and without flipping the
classification, and at $10^{-6}$ none. Second, while the integrator
dominates the error budget, the integer residual tracks the requested
tolerance with a roughly constant factor,
\begin{equation}\label{eq:rule}
  \varepsilon \;\approx\; 100 \cdot \mathrm{tol}
\end{equation}
to within an order of magnitude; the factor is inherited from the
accepted-step count of the march. At the standard window $\wmax=10^{4}$ the
truncated tail of this velocity-feedback system, $\approx2\cdot10^{-5}$,
takes over once the tolerance drops below $\sim10^{-6}$ -- the mean
residual of the tightest panel settles exactly there -- and, because every
error source surfaces in the same number, the diagnostic itself tells the
user that the next decade of certainty must come from $\wmax$, not from the
tolerance. The accuracy of $\Zint$ can therefore be \emph{prescribed}
rather than discovered: a user who wants the count certain to $10^{-3}$
asks for $10^{-5}$, which is exactly the package default. Third, the
CPU-time premium of tightening by four decades is a factor of three
(Table~\ref{tab:tolerance}), because the extra effort is spent only near
the integrand peaks.

For chart work a miscount is a milder error than a displaced boundary, and
the two degrade differently: the count is a rounded quantity and fails
discontinuously, whereas the tracked $\sest$ inherits the tolerance
smoothly. The traced boundary is built from $\sest$ alone through the
objective~\eqref{eq:objective}, and it stays a small fraction of one grid
pixel away from the reference trace at every tolerance studied
(\code{data/boundary\_shift.csv}). The practical recipe follows: a
moderately loose tolerance ($10^{-3}$--$10^{-4}$) is an appropriate tool
for an initial search -- the residual field will announce if it was too
loose -- after which the interesting region is recomputed at the default
setting to certify the counts, with the cross-check~\eqref{eq:crosscheck}
guarding the one error mode near the boundary that the residual cannot
reveal.

A flagged pixel need not be recomputed with a different backend: the march
can repair itself (\code{peak\_repair=true}). The skipped peak is easy to
\emph{locate} even when it is far too narrow to \emph{integrate}: its center
is a minimum of $|\Dfun(\sigma+\ii\omega)|^2$, which dips toward that
minimum over an $O(1)$ frequency range no matter how narrow the peak is. A
root-finding callback stops the integrator at every such minimum,
recomputes the phase increment of the crossing step by quadrature whose
subdivision is forced into the located peak, and restarts the march at the
minimum. Measured across the Hopf point of this chart
(\code{data/peak\_repair.csv}): the plain march loses one count from
$\Delta P = \RepairPlainFailDP$ inward, the Gauss--Kronrod backend -- which
must \emph{discover} the peak -- loses it silently at
$\Delta P = \RepairDeepDP$, whereas the repaired march, which \emph{knows}
where the peak is, still counts correctly there. The repair matters only in
the immediate vicinity of the boundary, and arming it for a whole sweep
costs $\approx\RepairCostFactor\times$ the plain march, which is why it is
off by default; the economical pattern is to sweep plainly, flag
with~\eqref{eq:crosscheck}, and re-evaluate only the flagged pixels with
the repair on.


\subsection{Rightmost roots and their refinement}\label{sec:showcase-roots}

\begin{figure}[tb]
  \centering
  \maybefig{fig_refinement}{\linewidth}
  \caption{The optional root refinement, judged by method-independent metrics
  on two systems -- a smooth quasi-polynomial (the showcase, top) and a
  rational frequency-response function with poles (the turning-lobe example of
  Appendix~\ref{app:gallery}, middle). Error \emph{fields}: per-pixel
  $\log_{10}|\Real\lam_{\mathrm{dom}}-\sigma_{\mathrm{SD}}|$, the deviation of
  the dominant root's real part from an independent semi-discretization
  reference, for the raw estimate and the four polishes. Error \emph{bars}
  (bottom): distributions, over the stable pixels, of the residual
  $|\Dfun(\lam_{\mathrm{dom}})|$ (distance to an actual root) and of the same
  $\sigma$-deviation, for each system, together with the per-point cost. The
  raw estimate is the smoothest field; every polish sharpens the smooth
  system uniformly, while on the rational system Newton becomes patchy (it can
  jump to a neighbouring root near a pole) and the Taylor and combined polishes
  stay consistent.}\label{fig:refinement}
\end{figure}

\begin{table}[tb]
  \centering
  \caption{Median error of the dominant root's decay rate $\Real\lam$ against
  a semi-discretization reference, for each refinement, on the smooth
  (showcase) and non-smooth (turning) systems of Fig.~\ref{fig:refinement}.
  Times are per parameter point for the complete solve; overhead is measured
  against the raw (unrefined) solve.}\label{tab:refinement}
  \maybeinput{tables/tab_refinement.tex}
\end{table}

Figure~\ref{fig:refinement} quantifies the rightmost-root estimates with a
yardstick that favours no method: the residual $|\Dfun(\lam)|$, which
measures distance to an actual root, and the deviation from an independent
semi-discretization solve.\footnote{A natural but circular way to rank the
polishes is to measure each against many Newton iterations taken as
``truth'', which trivially makes a few Newton steps look best; the metrics
used here avoid that trap.} Three facts follow. First, the raw tracked value
$(\sest,\west)$ already places the decay rate to $\RefShowRaw$ on the smooth
system and $\RefTurnRaw$ on the rational one, its sign -- all a chart needs
-- is correct throughout, and being a single closed-form expression per
pixel it is also the smoothest field, which is why it is the default.
Second, no single polish wins everywhere: on the smooth system Newton
drives the residual $|\Dfun|$ to machine precision and saturates the
accuracy of the semi-discretization reference ($\RefShowNewton$), but on
the rational system,
where the plant response has poles near the imaginary axis, a Newton step
can settle on a neighbouring root and its error \emph{worsens} to
$\RefTurnNewton$ -- the patchiness visible in its field -- whereas the cubic
Taylor polish stays consistent at $\RefTurnPolyThree$, and the combined mode
is the robust compromise ($\RefTurnComb$) that never speckles. Third, the
choice is nearly free: even the combined mode adds only about
$\RefCombOver\%$ to the complete solve (Table~\ref{tab:refinement}). The
practical hierarchy of Section~\ref{sec:refinement} follows directly: raw
for charts, a polish when a specific root matters, $\sigma$-contours
(Section~\ref{sec:showcase-sigma}) for a precise decay-rate map. Tracking
$N_r>1$ minima returns several dominant roots at once, which the multi-root
recipe of Section~\ref{sec:performance-sd} exploits.

\subsection{Decay-rate design: $\sigma$-contours}\label{sec:showcase-sigma}

\begin{figure}[tb]
  \centering
  \maybefig{fig_sigma_contours}{0.72\linewidth}
  \caption{Decay-rate map of the showcase system: colored background --
  estimated spectral gap $\sest$ from the coarse sweep (stable domain only);
  curves -- guaranteed $\gamma$-stability boundaries traced by MDBM for a
  family of shifted lines $\Real\lam=\sigma$. The two independent
  computations delineate the same landscape.}\label{fig:sigmacontours}
\end{figure}

Finally, Fig.~\ref{fig:sigmacontours} demonstrates the $\gamma$-stability
capability: the traced $\sigma$-contours (Section~\ref{sec:sigmacontours})
enclose progressively smaller regions with guaranteed spectral gaps, and
they are consistent with the interpolated $\sest$ background computed from
the coarse sweep alone. For control design this map answers directly which
gain combinations guarantee a prescribed settling rate -- at the cost of one
extra boundary tracing per contour.